\chapter{Satzbau}
\label{cha:Satzbau}

Der Satzbau im Deutschen folgt klaren Regeln, die die Wortstellung im Hauptsatz und Nebensatz bestimmen.

\section{Die Satzklammer}
\label{sec:Satzklammer}

Das deutsche Verb bildet eine „Satzklammer":

\begin{tcolorbox}[title=Satzklammer]
\begin{center}
Position 2: finites Verb \\
Ende: infinite Verbformen (Partizip, Infinitiv)
\end{center}
\end{tcolorbox}

\begin{itemize}
    \item Ich \textbf{sehe} den Film heute Abend \textbf{anzusehen}. (trennbares Verb)
    \item Er \textbf{hat} yesterday \textbf{arbeiten müssen}. (Modalverb)
    \item Wir \textbf{wollen} morgen ins Kino \textbf{gehen}.
\end{itemize}

\section{Hauptsatz-Wortstellung}
\label{sec:Hauptsatz}

\begin{center}
\begin{tabular}{|c|c|c|c|c|c|}
\hline
\textbf{Position 1} & \textbf{Position 2} & \textbf{Position 3} & \textbf{Position 4} & \textbf{Position 5} & \textbf{Position 6+} \\ \hline
Subjekt & finites Verb & Objekt 1 & Adverb & Objekt 2 & Verb 2... \\ \hline
\end{tabular}
\end{center}

Beispiele:
\begin{itemize}
    \item Ich \textbf{sehe} \underline{morgen} \underline{die} \underline{Freunde}.
    \item Heute \textbf{geht} \underline{er} \underline{nach} \underline{Hause}.
    \item \underline{Ich} \textbf{bin} \underline{nach} \underline{Frankfurt} \underline{gefahren}.
\end{itemize}

\section{Nebensatz-Wortstellung}
\label{sec:Nebensatz}

Im Nebensatz steht das Verb am Ende:

\begin{tcolorbox}[title=Nebensatz]
\begin{center}
Subjekt - Objekt - Adverb - \textbf{Verb am Ende}
\end{center}
\end{tcolorbox}

\begin{itemize}
    \item ...weil er \textbf{morgen} \underline{nach} \underline{Hause} \textbf{geht}.
    \item ...dass sie \underline{viel} \underline{Geld} \textbf{verdient}.
    \item ...obwohl es \textbf{kalt} \textbf{war}.
\end{itemize}

\section{Trennbare Verben}
\label{sec:Trennbar}

Im Hauptsatz: Finite Verbform Position 2, Partikel am Ende.
Im Nebensatz: Zusammen am Ende.

\begin{enumerate}
    \item Wann \underline{\hspace{1cm}} ihr in Berlin \underline{\hspace{1cm}}? (ankommen) \quad \textit{(Lösung: kommt ... an)}
    \item Um 8 Uhr \underline{\hspace{1cm}} der Unterricht \underline{\hspace{1cm}}. (anfangen) \quad \textit{(Lösung: fängt ... an)}
    \item Er \underline{\hspace{1cm}} mit dem Rauchen \underline{\hspace{1cm}}. (aufhören) \quad \textit{(Lösung: hört ... auf)}
    \item Wann \underline{\hspace{1cm}} du \underline{\hspace{1cm}}? (zurückkommen) \quad \textit{(Lösung: kommst ... zurück)}
    \item Ich \underline{\hspace{1cm}} ein Geschenk \underline{\hspace{1cm}}. (mitbringen) \quad \textit{(Lösung: bringe ... mit)}
    \item Das Kind \underline{\hspace{1cm}} um 20 Uhr \underline{\hspace{1cm}}. (einschlafen) \quad \textit{(Lösung: schläft ... ein)}
    \item Wann \underline{\hspace{1cm}} ihr \underline{\hspace{1cm}}? (vorbeikommen) \quad \textit{(Lösung: kommt ... vorbei)}
    \item Sie \underline{\hspace{1cm}} nach der Arbeit \underline{\hspace{1cm}}. (weggehen) \quad \textit{(Lösung: geht ... weg)}
    \item Ich \underline{\hspace{1cm}} dich später \underline{\hspace{1cm}}. (anrufen) \quad \textit{(Lösung: rufe ... an)}
    \item Kannst du das Fenster \underline{\hspace{1cm}}? (aufmachen) \quad \textit{(Lösung: aufmachen)}
\end{enumerate}

\chapterend